% Supplementary material, rebuilt on 2026-11-16 (layout pass). Of the former
% Appendices A to C, only what the chapters need for reproducibility stays
% here; everything else is in docs/thesis/defence-notes.md, verbatim. The
% lineage table carries the two corrections of the Pipeline History Workbench.
\chapter{Supplementary Material}
\label{app:detail}
\label{app:figures}

This appendix holds the pipeline diagram and the lineage of \cref{ch:content},
the parameters and the seed sweep of the study with scripted readers of
\cref{ch:loop}, the database rules that \cref{ch:loop,ch:trust} rely on, and
who built which part of the code (\cref{ch:introduction}).

The code at the pin is public in the repository
\url{https://github.com/osamash0/ascend-academy}, revision
\code{dfc20df61d0caf8fb1c9\allowbreak 0a215fa72bcd394de4a1}. The evidence that the
chapters cite under \code{docs/\allowbreak thesis/} holds no slide text and
no rows of the production database. It is part of the digital submission of
this thesis, with every commit that the chapters name, as a copy of the
repository (a git bundle, branch \code{thesis-evidence}).

\section{The Pipeline and Its Lineage (Chapter 4)}
\label{sec:lineage-table}

\Cref{tab:lineage} lists by month the pipeline versions v1 to v5 of
\cref{sec:c4-lineage}, which led to the pipeline of \cref{sec:pipeline}; the
filter of v1 and three later rebuilds of v5 have rows of their own.

\begin{table}[tp]
  \centering
  \small
  \caption{The pipeline lineage at a glance (2026; v1 to v5 are the numbered
    pipeline versions, unrelated to the v4 interface of \cref{sec:vision}).
    Sources: the commits cited in \cref{sec:c4-lineage} and, for the last
    three rows, PRs~\#84, \#139, \#145, and \#178; the rows v1 and v5 were
    checked commit by commit in November 2026.}
  \label{tab:lineage}
  \begin{tabular}{@{}llp{0.6\textwidth}@{}}
    \toprule
    Generation & Month & What it did, and what became of it \\
    \midrule
    v1 slide enhancement & Feb & Model calls per page at upload: three
      (rewrite, summary, quiz), four from March with a title call, one from
      late March as a single structured call, and from late April one vision
      call per page when a vision model was available; the original page
      shown next to the AI text. Replaced in May. \\
    content filter & Mar & Pattern check, density score, then a model
      judge: non-teaching pages get no AI content. No longer called after
      June. \\
    v2 routed pipeline & May & Each page classified, then sent to a text,
      vision, or table route (tables by OCR first); a deck-level ``blueprint''. The default until June. \\
    v3 staged pipeline & May & Five resumable stages (eight in the design),
      Docling, an Arq job queue, pgvector, a LiteLLM gateway. Never the
      default; later described in the code as dormant. \\
    v4 orchestrator & May & A single model stage shared by all external
      extractors. In an offline routing check in June, 142 of 143 pages of
      one deck went to text. \\
    v5 unified & Jun & One server-side orchestrator: the run table and job
      queue of v3, the prompts of v4 in one call per page, and a vision model
      for pages with fewer than 25 characters of text (the engine of v2 only
      in its first commit); the server writes the result. The only reachable
      path since July. \\
    deck plan first & Oct & Plan the deck (roles, sections, boilerplate) with
      no model call, then generate; parse receipts with cost. In use. \\
    page by page & Oct & The reading view opens at once and pages arrive as
      they are written; a re-saved PDF copy for the reading view. In use. \\
    one routing rule & Nov & Text, figure, or no-text decided in one place
      from the page itself. In the code at the pin; deployed since November
      2026. \\
    \bottomrule
  \end{tabular}
\end{table}


\section{The Scripted Readers in Detail (Chapter 5)}
\label{sec:c5-personas-detail}

The study of \cref{sec:c5-evaluation} builds, for each arm, a Space with the same Lesson of
30~pages. Five ideas are anchored on pages 5, 10, 15, 20, and 25, where a
reader checks in and meets a checkpoint, and pages 8 to 12 are the hard
pages. Each member reads the Lesson once, on a day drawn uniformly from the
55~days before a fixed clock (from the six days before it in arm~I). Arm
number $k$ of \cref{tab:persona}, counted from zero, runs with seed
$20261109 + k$. We chose every behaviour and value in
\cref{tab:personas-params} ourselves; none comes from data, and the askers
draw their questions from six that we wrote. Two personas send a third
check-in, \emph{unsure}, that the reading view no longer offers
(\code{v4\_reader.py:340}). It counts like \emph{Clicked} in the share of
\emph{Still foggy} check-ins, and like \emph{Still foggy} in the share of
readers who reached for help and in the calibration gap. The script is archived next to
its results (\code{docs/\allowbreak thesis/\allowbreak experiments/\allowbreak
persona-study/\allowbreak persona\_study.py}).

\begin{table}[htbp]
  \centering
  \small
  \caption{The personas. A reading time is drawn from a normal distribution
    with the given mean and standard deviation and cut to the range in
    brackets.}
  \label{tab:personas-params}
  \newcommand*{\rr}{\raggedright\hspace{0pt}}
  \begin{tabular}{@{}lp{0.22\textwidth}p{0.15\textwidth}p{0.19\textwidth}p{0.14\textwidth}@{}}
    \toprule
    Persona & \rr Reading time per page & \rr Explain, Ask & \rr Check-in & \rr Checkpoint \tabularnewline
    \midrule
    steady & \rr 40\,s, sd 10 (20--60), every page & \rr none & \rr \emph{Clicked} & \rr right \tabularnewline
    struggling & \rr as steady, three times as long on pages 8--12; turns back from page~13 to page~12 twice & \rr Explain on pages 8--12, Ask there with probability 0.7 & \rr \emph{Still foggy} on page~10, elsewhere \emph{Clicked}, \emph{Clicked}, or \emph{unsure} at random & \rr wrong on page~10, elsewhere right with probability 0.8 \tabularnewline
    asker & \rr 40\,s, sd 10 (20--60) & \rr Ask on every third page & \rr \emph{unsure} on page~10, elsewhere \emph{Clicked} & \rr right with probability 0.85 \tabularnewline
    skimmer & \rr 2.5\,s, sd 1 (1--4.5), pages 1--10 only & \rr none & \rr none & \rr none \tabularnewline
    overconfident & \rr 35\,s, sd 8 (20--60) & \rr none & \rr \emph{Clicked} & \rr wrong on page~10, elsewhere right \tabularnewline
    confused asker\textsuperscript{a} & \rr 40\,s, sd 10 (20--60) & \rr Ask on pages 8--12 & \rr \emph{Still foggy} on page~10, elsewhere \emph{Clicked} & \rr wrong on page~10, elsewhere right \tabularnewline
    \bottomrule
  \end{tabular}\par\smallskip
  {\footnotesize \textsuperscript{a}Post hoc, arm~J only.}
\end{table}

After the result we ran all ten arms with 100 seed bases, under a protocol
committed before the run (\code{seed-sweep/\allowbreak PROTOCOL.md}, commit
\code{1db7b660}). Base~0 is the published run, and the sweep reproduced it.
\Cref{tab:sweep} gives, for each expectation written before the first run,
the number of seeds in which it held. The predictions in the protocol held
except two, narrowly: we expected the hotspot of arm~A in 20 to 60 of 100
seeds and saw it in 19, and we expected no calibration gap in arm~A and saw
one, which \cref{sec:c5-results} traces to chance.
Of the post-hoc expectations, the hotspot of arm~J held in all 100
seeds, its themes in 61, and the numbers of arm~I in all 100.

\begin{table}[htbp]
  \centering
  \small
  \caption{The 23 expectations written before the first run, over 100 seed
    bases (\code{seed-sweep/\allowbreak RESULTS.md}). Per seed the engine met
    16 to 19 of them (median 17); an engine that reports nothing meets 16.}
  \label{tab:sweep}
  \begin{tabular}{@{}llr@{}}
    \toprule
    Arm & Expectation & Held in \\
    \midrule
    A & a hotspot & 19 of 100 \\
    A & no skipped page; no calibration gap & 100; 99 of 100 \\
    A & the weekly numbers have a value & 1 of 100 \\
    A & the reach word is not \code{few} & 100 of 100 \\
    B & no hotspot; no finding on pages 8--12 & 0; 0 of 100 \\
    C & themes are made & 0 of 100 \\
    D & no themes & 100 of 100 \\
    E & a leak; skipped pages; no hotspot & 100 of 100 each \\
    F & a calibration gap & 100 of 100 \\
    G & nothing shown, numbers empty, reach \code{few} (nine expectations) & 100 of 100 each \\
    H & the weekly numbers have a value & 17 of 100 \\
    \bottomrule
  \end{tabular}
\end{table}

\section[Declared and Deployed Database Rules (Chapters 5 and 6)]{Declared and Deployed Database Rules \mbox{(Chapters 5 and 6)}}
\label{sec:deployed-rules}

A read-only probe of the production catalogue in November 2026 read the
policies and function definitions that \cref{tab:deployed} lists. Its entries
are readings of what each rule permits; we tested no exploit.

\begin{table}[H]
  \centering\small
  \caption{Declared, tested, and deployed: database-side rules that bear on
    \cref{ch:loop,ch:trust}. The middle column is what the migrations at the
    pin declare (database tests exist for the partitioning, the role at sign-up, the
    professor policy, and the anonymous SECURITY DEFINER exposure), the right
    column what the read-only probe of production found in November 2026.}
  \label{tab:deployed}
  \begin{tabular}{@{}p{0.27\textwidth}p{0.3\textwidth}p{0.35\textwidth}@{}}
    \toprule
    Rule & Declared (migrations at the pin) & Deployed (November 2026) \\
    \midrule
    Event table partitioned by month &
      declared; the tests check the partitions &
      plain table, no partitions \\
    Role at sign-up &
      student or professor, as requested &
      student only; professor set by an administrator \\
    Professor reads members' events &
      policy present &
      present; covers a member's events in every Space \\
    Lecture rows: course ownership &
      not checked by the insert policy; checked in the API &
      not checked by the insert policy \\
    Self-enrolment &
      a student may enrol in any course &
      as in the repository; publication not checked \\
    Lecture posters (page~1 image) &
      readable by every signed-in account &
      as in the repository, private uploads included, given the path \\
    \code{get\_user\_profile} &
      executable by every signed-in account &
      as in the repository; returns name, picture, and institution for any id \\
    \code{SECURITY DEFINER} functions &
      anonymous execution revoked (October) &
      49 executable by \code{authenticated}; whether each checks its caller was not audited \\
    \bottomrule
  \end{tabular}
\end{table}

\section{Who Built What (Chapter 1)}
\label{sec:contributions}

\begin{table}[H]
  \centering\footnotesize
  \renewcommand*{\code}[1]{\texttt{#1}}%
  \caption{Lines of code at the pin by the developer whose commit last changed
    them (\code{git blame -w -C -M}), without tests and generated, lock, and
    data files. A: the author, with the coding agents the author directed; C:
    the co-developer. Of the 1{,}528 commits that the pin reaches, 53 are C's;
    eight more come from a bot identity, whose 65 lines the author counts for
    C. All evaluation code is A's: \code{backend/eval/} at the pin and every
    commit after the pin. C created the upload service and its validation (\code{21b4d8a2},
    \code{377f7fa1}); C's security fixes of July 2026 survive in 359 lines, the
    extractor wiring (\code{8aa255e2}) in 36, and the Groq integration
    (\code{76aa4b42}) in 4.}
  \label{tab:contrib}
  \begin{tabular}{@{}p{0.62\textwidth}rrl@{}}
    \toprule
    Part, with its main paths at the pin & A & C & Created in \\
    \midrule
    Pipeline (RQ1): \code{parser/}, \code{arq\_worker.py} & 4{,}533 & 26 & A, \code{6ce7b5f1} \\
    Upload entry (RQ1): \code{upload.py}, \code{upload\_service.py}, validation & 1{,}192 & 214 & A, \code{be25cc1d} \\
    Benchmark (RQ1): \code{eval/parser\_bench/} & 9{,}518 & 0 & A, \code{e2cc63c5} \\
    Events, engine, owner's view (RQ2): \code{insights/}, \code{landing/} & 5{,}325 & 0 & A, \code{0b2aa58e} \\
    Quiz checks (RQ3): \code{quiz\_grounding.py}, answer check & 614 & 0 & A, \code{323e234b} \\
    Reading interface: \code{src/features/spaces/}, v4 routes & 43{,}337 & 0 & A, \code{f9aad62c} \\
    \midrule
    Database schema: \code{supabase/migrations/}, other than events & 10{,}126 & 355 & C, \code{184e5742} \\
    Stack: \code{package.json}, build configs, \code{App.tsx}, \code{main.py} & 2{,}682 & 641 & C, \code{184e5742} \\
    Sign-in: \code{auth.tsx}, \code{Auth.tsx}, \code{auth\_middleware.py} & 1{,}122 & 200 & C, \code{184e5742} \\
    Content filter: \code{content\_filter.py} & 60 & 315 & C, \code{bfbe6c09} \\
    Lecturer analytics: \code{analytics\_service.py} and its page & 1{,}763 & 103 & C, \code{8a34f6a0} \\
    Admin dashboard: \code{AdminDashboard.tsx}, \code{admin.py}, \code{admin/} & 1{,}040 & 1{,}265 & A, \code{343890c5} \\
    Tutor memory: \code{chat\_memory.py}, two services in \code{ai/} & 286 & 462 & C, \code{c96a4d32} \\
    \midrule
    All code at the pin, 1{,}148 files & 191{,}969 & 6{,}887 & \\
    \bottomrule
  \end{tabular}
\end{table}

\thesisfigurepage{f8-pipeline}
  {The ingestion pipeline. Two entry routes admit a PDF synchronously and
   converge on one content-addressed object and one Arq job, which writes each
   page as soon as it is synthesised. Pentagons mark where the reader gains
   something --- the original PDF at once, then each persisted page.}
  {fig:pipeline}
